\section{SMT-Based Constraint Search and the 18-Vertex Depth-2 Discovery}
\label{sec:discussion}

The near-path sweeps for $m \le 11$ and the 15-vertex bipartite example
$H^*$ illustrate how quickly the Leontovich search space grows. Beyond the
range of exhaustive enumeration, we encode the homomorphism inequalities as
an integer constraint problem for a Satisfiability Modulo Theories (SMT)
solver. This search produces the 18-vertex depth-2 bipartite example
$H_{18}$ studied below.

\subsection{SMT-based algebraic search}

To bypass the exponential explosion of graph enumeration (e.g., more than $5 \times 10^{11}$ bipartite targets on $m = 16$ vertices),
the search for a Leontovich target graph $H$ of order $m$ can be formulated as a Constraint Satisfaction Problem (CSP)
over the binary adjacency variables $A_{ij} \in \{0, 1\}$.

Rather than explicitly unrolling the sum over all $m^n$ possible mappings to compute tree homomorphisms,
we embed the dynamic programming walk recurrences directly into the solver's constraint network as auxiliary variables.
For a target graph $H$, we define the walk vectors $w_k \in \mathbb{Z}^m$ recursively:
\begin{align}
  w_0(i) &= 1, \quad \forall i \in \{1, \dots, m\} \\
  w_k(i) &= \sum_{j=1}^{m} A_{ij} \, w_{k-1}(j), \quad \forall k \ge 1
\end{align}
The homomorphism count for the path is given by the linear sum
$\Hom(P_n, H) = \sum_{i=1}^{m} w_{n-1}(i)$. For the near-path family
$E_n^{(d)}$ with a pendant edge attached at depth $d$, the count can be
represented compactly as:
\begin{equation}
  \Hom(E_n^{(d)}, H) = \sum_{i=1}^{m} w_{n-d-2}(i) \, w_1(i) \, w_d(i)
\end{equation}
The search for a Leontovich target graph then reduces to solving the polynomial inequality:
\begin{equation}
  \label{eq:smt_ineq}
  \Hom(E_n^{(d)}, H) < \Hom(P_n, H)
\end{equation}
subject to adjacency constraints. Bipartiteness can be enforced dynamically
by introducing partition variables $x_i \in \{0, 1\}$ such that $A_{ij} = 1
\implies x_i \neq x_j$, while connectivity is guaranteed by defining
distance levels from a root vertex.

\subsection{Local SMT Leaf-Pruning Proof for \texorpdfstring{$T(7,1,9)$}{T(7,1,9)}}
\label{sec:local-pruning-audit}

To evaluate the optimality of our 76-vertex Leontovich graph
$T(7,1,[9^5, 8^2])$ discovered by simulated annealing, we formulated a
highly targeted local leaf-pruning search. Starting from the fully symmetric
tree $T(7,1,9)$ (78 vertices), we define integer decision variables $k_1,
\dots, k_7 \in \{0, \dots, 9\}$ representing the number of active leaves on
each of the 7 hubs. To break isomorphic permutations, we enforce $9 \ge k_1 \ge
k_2 \ge \dots \ge k_7 \ge 0$.

Because of the high symmetry of the quotient orbits, the walk equations can be
simplified from $m = 15 + \sum k_i$ variables down to a bilinear recurrence of
exactly $22$ variables per step:
\begin{align}
  w_k(\text{root}) &= \sum_{i=1}^{7} w_{k-1}(\text{bridge } i) \\
  w_k(\text{bridge } i) &= w_{k-1}(\text{root}) + w_{k-1}(\text{hub } i) \\
  w_k(\text{hub } i) &= w_{k-1}(\text{bridge } i) + k_i \, w_{k-1}(\text{leaf-type } i) \\
  w_k(\text{leaf-type } i) &= w_{k-1}(\text{hub } i)
\end{align}
We executed this local SMT solver to check if any pruning of size $|V| \le 75$ is
Leontovich for any threshold $n \le 31$. The solver returned \textbf{UNSAT} for every
single threshold.

This establishes a bounded computational audit:
\begin{compresult}
  Within the $T(7,1,9)$ star-of-stars pruning landscape, no pruning with
  $|V| \le 75$ is detected as Leontovich by the completed integer constraint
  checks for thresholds $n \le 31$. The 76-vertex tree
  $T(7,1,[9^5, 8^2])$ is therefore the minimum witness certified by this
  bounded audit.
\end{compresult}

\subsection{The 18-Vertex Depth-2 Bipartite Leontovich Graph Discovery}

To search for small non-tree Leontovich graphs, we restricted the
target class to bipartite graphs with left-partition size $m_1 \in \{2, 3\}$.
Under this partition constraint, every vertex in the right partition $R$ must
connect to a non-empty subset of the left partition $L$.

For $m_1 = 2$, there are only $2^2 - 1 = 3$ valid connection patterns for
vertices in $R$. A graph is thus uniquely specified up to isomorphism by a
triplet $(c_1, c_2, c_3)$ where $c_1 + c_2 + c_3 = m_2$.

\begin{proposition}[The $n=5$ depth-2 obstruction for $m_1 = 2$]
For every connected bipartite graph of partition $(2,m_2)$, the depth-2
near-path comparison has no crossover at $n=5$.
\end{proposition}
\begin{proof}
Let $H$ be a connected bipartite graph with left partition $L = \{u_1, u_2\}$.
The homomorphism counts at odd thresholds $n=2k+1$ are governed by the $2 \times 2$ transfer matrix $M = A^2|_L$, where $A$ is the adjacency matrix. We have:
\[
  M = \begin{pmatrix} c_1+c_3 & c_3 \\ c_3 & c_2+c_3 \end{pmatrix}.
\]
For $H$ to be connected, we must have $c_3 \ge 1$. Direct algebraic expansion of the transfer matrix product for $n=5$ yields the exact difference:
\[
  \Hom(P_5, H) - \Hom(E_5^{(2)}, H) = -c_3 \bigl( (c_1 - c_2)^2 + c_1 + c_2 \bigr).
\]
Since $c_1, c_2 \ge 0$ and $c_3 \ge 1$, this difference is non-positive, meaning $\Hom(P_5, H) \le \Hom(E_5^{(2)}, H)$.

\end{proof}

Beyond $n=5$, every $(2,m_2)$ target of total order $m\le15$ is covered at
every tested depth by the exhaustive bipartite sweep of
Section~\ref{sec:computation}, which found no crossover for odd $n\le199$ and
$d\le20$. We make no unconditional claim for larger $m_2$ or beyond that
window.

For $m_1 = 3$, there are $2^3 - 1 = 7$ valid connection patterns: 3
singletons ($c_1, c_2, c_3$), 3 pairs ($c_4, c_5, c_6$), and 1 triplet ($c_7$).
Any such bipartite target is uniquely determined by the vector $(c_1, \dots,
c_7)$.

The 15-vertex graph $H^*$ from Remark~\ref{rem:depth_anomaly} shows that
allowing arbitrary near-path depths yields a smaller bipartite Leontovich graph.
The following result isolates the smallest example found by the stricter
depth-2 test $E_n^{(2)}$, the near-path family used in the classical
Leontovich construction:
\begin{compresult}
  The bipartite graph $H_{18}$ of partition $(3, 15)$ with pattern vector
  $(7, 0, 0, 1, 1, 6, 0)$ is the unique minimal depth-2 bipartite Leontovich
  graph of partition $(3, m_2)$ found in this completed depth-2 bounded sweep
  through path length $n \le 51$. It first crosses against
  $E_n^{(2)}$ at $n = 17$. Its Graph6 representation is
  \begin{equation*}
    \texttt{QCaCCA?\_E?S?W?W?K?B??W?@\_??}
  \end{equation*}
\end{compresult}

Allowing arbitrary attachment depth gives an earlier crossover for the same
target. Exact integer arithmetic gives
\begin{align*}
  \Hom(P_{15},H_{18})&=1{,}105{,}256{,}666,\\
  \Hom(E_{15}^{(4)},H_{18})&=1{,}104{,}756{,}090,
\end{align*}
so the first near-path crossover for $H_{18}$ occurs at $(n,d)=(15,4)$,
with margin $500{,}576$. The symmetric depth $d=9$ gives the same source tree,
since $E_{15}^{(4)}\cong E_{15}^{(9)}$ by Remark~\ref{rem:d-symmetry}.
An exact scan over every admissible depth finds no crossover for $n<15$.

At the depth-2 crossover threshold $n=17$, the homomorphism counts verified by exact-precision arithmetic are:
\begin{align*}
  \Hom(P_{17}, H_{18}) &= 14{,}801{,}051{,}732 \\
  \Hom(E_{17}^{(2)}, H_{18}) &= 14{,}795{,}982{,}954
\end{align*}
yielding a positive Leontovich margin of $\Delta = 5{,}068{,}778$. Since
$m=18 < 76$, this is less than one quarter the order of the smallest
previously known depth-2 construction in the tested search regime. Because
the exhaustive check of all 45,838
non-isomorphic configurations for $m_2 < 15$ returned zero depth-2 hits,
\textbf{18 vertices is the depth-2 minimum within this partition family in the
completed depth-2 sweep through path length $n \le 51$}.

\subsubsection{\texorpdfstring{Structure and walk iteration of $H_{18}$}{Structure and walk iteration of H18}}

\begin{figure}[H]
  \centering
  \includegraphics[width=0.9\textwidth]{figures/leontovich_18.pdf}
  \caption{\texorpdfstring{The optimal 18-vertex depth-2 bipartite Leontovich graph $H_{18}$ discovered via SMT synthesis.}{The optimal 18-vertex depth-2 bipartite Leontovich graph H18 discovered via SMT synthesis.}}
  \label{fig:walk_h18}
\end{figure}

To make the mechanics of the Leontovich crossover transparent, we detail the
orbit decomposition and walk recurrences for $H_{18}$.
By partitioning the $18$ vertices into canonical symmetry groups, we reduce the
$18 \times 18$ transition adjacency representation into a $7 \times 7$ quotient system. The seven symmetry groups are:
\begin{enumerate}
  \item $L_0$: Left vertex $v_0$ ($1$ node, degree $9$).
  \item $L_1$: Left vertex $v_1$ ($1$ node, degree $7$).
  \item $L_2$: Left vertex $v_2$ ($1$ node, degree $7$).
  \item $R_{\mathrm{leaf},0}$: Right leaves connected only to $v_0$ ($7$ nodes, degree $1$).
  \item $R_{01}$: Cross-connector connected to $\{v_0, v_1\}$ ($1$ node, degree $2$).
  \item $R_{02}$: Cross-connector connected to $\{v_0, v_2\}$ ($1$ node, degree $2$).
  \item $R_{12}$: Bipartite core connected to $\{v_1, v_2\}$ ($6$ nodes, degree $2$).
\end{enumerate}

The step-by-step walk recurrence vectors $w_k(u)$ (denoting the number of walks of length $k$ starting at $u$)
are governed by the following exact coupled system:
\begin{align*}
  w_k(L_0) &= 7 w_{k-1}(R_{\mathrm{leaf},0}) + w_{k-1}(R_{01}) + w_{k-1}(R_{02}) \\
  w_k(L_1) &= w_{k-1}(R_{01}) + 6 w_{k-1}(R_{12}) \\
  w_k(L_2) &= w_{k-1}(R_{02}) + 6 w_{k-1}(R_{12}) \\
  w_k(R_{\mathrm{leaf},0}) &= w_{k-1}(L_0) \\
  w_k(R_{01}) &= w_{k-1}(L_0) + w_{k-1}(L_1) \\
  w_k(R_{02}) &= w_{k-1}(L_0) + w_{k-1}(L_2) \\
  w_k(R_{12}) &= w_{k-1}(L_1) + w_{k-1}(L_2)
\end{align*}
with the base canonical condition $w_0(u) = 1$ for all vertices. The explicit numerical values for the
first five steps are shown in Table~\ref{tab:walk_h18}. This iteration can be carried out exactly to $w_{16}(u)$.

\begin{table}[H]
\centering
\caption{Step-by-step walk vector iteration and total sum counts for $H_{18}$.}
\label{tab:walk_h18}
\begin{tabular}{lcccccccc}
  \toprule
  Step $k$ & $L_0$ & $L_1$ & $L_2$ & $R_{\mathrm{leaf},0}$ ($\times 7$) & $R_{01}$ ($\times 1$) & $R_{02}$ ($\times 1$) & $R_{12}$ ($\times 6$) & Total $\sum w_k$ \\
  \midrule
  $w_0$ & 1 & 1 & 1 & 1 & 1 & 1 & 1 & 18 \\
  $w_1$ & 9 & 7 & 7 & 1 & 2 & 2 & 2 & 46 \\
  $w_2$ & 11 & 14 & 14 & 9 & 16 & 16 & 14 & 218 \\
  $w_3$ & 95 & 100 & 100 & 11 & 25 & 25 & 28 & 590 \\
  $w_4$ & 127 & 193 & 193 & 95 & 195 & 195 & 200 & 2{,}768 \\
  $w_5$ & 1055 & 1395 & 1395 & 127 & 320 & 320 & 386 & 7{,}690 \\
  \bottomrule
\end{tabular}
\end{table}

The path count is given by the sum $\Hom(P_{17}, H_{18}) = \sum_u w_{16}(u)$, while the near-path count
is computed as $\Hom(E_{17}^{(2)}, H_{18}) = \sum_u w_{13}(u) w_1(u) w_2(u)$, yielding the exact margin reported above.

\begin{remark}[Rigorous Resolution of Other Bipartite Families]
  To determine whether smaller depth-2 examples occur in other bipartite
  partition families, we engineered a high-performance C++ solver to
  exhaustively sweep all other partition families $(m_1, m_2)$ for
  $m_1 \ge 4$ up to total order $m_1 + m_2 \le 17$. By evaluating all
  non-isomorphic bipartite connection patterns up to path length $n = 51$
  (covering millions of distinct configurations), the solver returned zero
  depth-2 hits. This supports the claim that $H_{18}$ is the smallest
  currently verified depth-2 bipartite Leontovich graph across the tested
  partition families in the completed depth-2 sweep through path length
  $n \le 51$; beyond that bounded regime, Section~\ref{sec:doublecover}
  shows why larger thresholds can behave differently. The depth-dependent
  graph $H^*$ remains the smaller bipartite Leontovich example for
  general $E_n^{(d)}$.
\end{remark}

\subsection{The Near-Miss Asymptotic Dilution}

During our trace of the $m_1 = 3$ family, we discovered pattern vectors where
the homomorphism count ratio approached $1.0$ from above with extreme
precision (e.g., the 14-vertex target $c = (9, 1, 0, 0, 0, 1, 0)$).
However, algebraic inspection reveals these are not impending crossovers, but
disjoint unions.

The vector $(9, 1, 0, 0, 0, 1, 0)$ specifies a disconnected graph
$H \cong K_{1,9} \sqcup P_4$. Because homomorphism counts map connected trees
entirely into a single connected component, they are strictly additive over
disjoint components: $\Hom(T, H) = \Hom(T, K_{1,9}) + \Hom(T, P_4)$. The star
$K_{1,9}$ is degenerate (Proposition~\ref{prop:p3}) and yields
identical counts for $P_n$ and $E_n^{(2)}$, contributing a massive
$O(3^n)$ perfect tie term. The path $P_4$ provides a strict positive margin,
but at a much slower growth rate $O(1.618^n)$. The ratio approaches $1.0$
precisely because it is physically impossible to cross over; the strict minimality of $P_4$
is asymptotically diluted by the dominant degenerate star.

\subsection{The parameterized \texorpdfstring{$\mathcal{B}(x,y)$}{B(x,y)} family}

The pattern $(7,0,0,1,1,6,0)$ of $H_{18}$ belongs to a natural
two-parameter family of bipartite targets.

Let $\mathcal{B}(x,y)$ be the bipartite graph defined by the left partition
$L = \{u_0, u_1, u_2\}$ and the connection pattern $(x, 0, 0, 1, 1, y, 0)$.
This graph has left partition size $m_1 = 3$ and right partition size
$m_2 = x + y + 2$, giving a total order of $m = x + y + 5$.

Due to the symmetry of the connection subsets, $u_1$ and $u_2$ form a perfect
automorphism orbit. The random walks restricted to the left partition taking
steps of length 2 are governed by a $3 \times 3$ matrix. However, the
antisymmetric eigenvector $(0, 1, -1)^T$ has an exact eigenvalue of $\lambda^2=1$,
but is perfectly orthogonal to the all-ones initial state vector $\mathbf{1}$.
Thus, it contributes nothing to the homomorphism counts. The active dynamics
collapse into the $2 \times 2$ asymmetric quotient matrix:
\begin{equation}
  M = A^2|_L = \begin{pmatrix} x+2 & 2 \\ 1 & 2y+1 \end{pmatrix}.
\end{equation}

The non-zero squared eigenvalues of $\mathcal{B}(x,y)$ are precisely the roots
of the characteristic polynomial of $M$:
\begin{equation}
  \lambda^2 - (x + 2y + 3)\lambda + (2xy + x + 4y) = 0.
\end{equation}

For our record-breaking graph $H_{18}$, we evaluate $\mathcal{B}(7,6)$. The
characteristic polynomial simplifies to $\lambda^2 - 22\lambda + 115 = 0$,
yielding exact squared eigenvalues of $11 \pm \sqrt{6}$. The active eigenvalue
ratio governing the crossover is thus exactly:
\begin{equation}
  r = \sqrt{\frac{11 - \sqrt{6}}{11 + \sqrt{6}}} \approx 0.797.
\end{equation}

The discriminant $\Delta = (x-2y+1)^2 + 8$ shows that the
ratio is maximized precisely along the ``resonance line'' $x = 2y - 1$,
which favors long-lived eigenvalue interference. Computations suggest that
this parameterized family contains infinitely many Leontovich graphs, but we
do not prove that assertion here. Within the completed
depth-2 $(3,m_2)$ block-pattern sweep through path length $n \le 51$,
$H_{18}$ is the smallest member of $\mathcal{B}(x,y)$ found to cross over.

% ════════════════════════════════════════════════════════════════════════
